\documentclass[12pt]{article}
\usepackage{sectsty,setspace}
\usepackage[top=1.87cm, bottom=1.87cm, left=1.87cm, right=1.87cm]{geometry} 
\usepackage{epstopdf}
\usepackage{amsmath,latexsym,amssymb,wasysym}
\usepackage{times}
\usepackage{fancyhdr}
\usepackage[T1]{fontenc}
\usepackage[skip=10pt plus1pt, indent=0pt]{parskip} % ligne vide entre paragraphes = espace automatique, plus besoin de \\

\setstretch{1}
%\setlength{\baselineskip}{0.167in}
\setlength\parindent{0pt}

\pagestyle{fancy}
\fancyhf{}
\fancyhead[R]{Christophe Rouleau-Desrochers}
\fancyfoot[C]{\thepage}
\renewcommand{\headrulewidth}{0pt}

\begin{document}
% ==============================================================================
\subsection*{I --- Expériences pertinentes}
% ==============================================================================
% <><><><><><><><><><><><><><><><><><><><><><><><><><><><><><><><><><><><><><><>
\textbf{a. Expériences de recherche} 
% <><><><><><><><><><><><><><><><><><><><><><><><><><><><><><><><><><><><><><><>

% --- --- --- --- --- --- --- --- --- --- --- --- --- --- --- --- --- --- ---
% Item 1 %
% --- --- --- --- --- --- --- --- --- --- --- --- --- --- --- --- --- --- ---
\textbf{Maîtrise en sciences forestières, University of British Columbia} --- sept. 2024 à aujourd'hui. Direction~: E.M. Wolkovich. \textit{Projets de recherche intégré au cursus.} 

J'ai conçu et dirigé une expérience en serre sur sept espèces portant sur la longueur et les conditions des saisons de croissance~: j'en ai formulé les hypothèses, défini les traitements et intégré les données climatiques et phénologiques à mes modèles statistiques. Au cours de cette expérience, j'ai formé et supervisé 15 étudiantes et étudiants de premier cycle aux protocoles de chambre de croissance, au suivi phénologique et à la gestion de serre. La suite de ce project deviendra le premier chapitre de mon doctorat. 

Le premier chapitre de ma thèse porte aussi sur la croissance de quatre espèces de Betulaceae issues de quatre provenances, en combinant la des anneux de croissance (échantillons d'arbres que j'ai sablés, scannés et mesurés), phénologie de début et de fin de saison pendant trois ans. Finalement, j'ai analysé les données avec l'aide de modèles bayésiens hiérarchiques. L'article est actuellement en révision chez \textit{Global Change Biology et a été soumis le 15 Septembre 2026}.

Finalement, le deuxième chapitre de ma thèse porte sur 11 espèces d'arbres matures et comment leur croissance est affectée par la longeur de leur saison de croissance et l'influence des conditions des années précédentes avec une perspective sur les traits d'histoire de vie et l'anatomie du bois. J'ai collecté les échantillons d'anneaux de croissance à Boston au printemps 2025, voyage durant lequel j'ai aussi présenté mes travaux à deux reprises à la communauté de science citoyenne Treespotters (2025--2026) et je garde contact avec eux régulièrement. Cet article sera soumis pour publication en octobre 2026. 

% --- --- --- --- --- --- --- --- --- --- --- --- --- --- --- --- --- --- ---
% Item 3 %
% --- --- --- --- --- --- --- --- --- --- --- --- --- --- --- --- --- --- ---
\textbf{Campagnes de terrain, Smithers (C.-B.) et Mount Rainier (Wash.)} --- étés 2024 et 2025. Direction~: E.M. Wolkovich. 

J'ai participé à des campagnes sur la dynamique de régénération forestière, où j'ai perfectionné mes méthodes forestières et appris les bases de la dendrochronologie. En 2025, j'ai dirigé l'équipe de terrain de Smithers en 2025~: planification des déplacements, répartition des tâches et contrôle de la qualité des données. 

% --- --- --- --- --- --- --- --- --- --- --- --- --- --- --- --- --- --- ---
% Item 4 %
% --- --- --- --- --- --- --- --- --- --- --- --- --- --- --- --- --- --- ---
\textbf{Bourse de recherche de premier cycle (BRPC/USRA), University of British Columbia} --- mai à août 2023. Direction~: Drs. F. Baumgarten et E.M. Wolkovich.

Premier contact avec les dispositifs expérimentaux de grande envergure~: j'ai conçu et mis en place une expérience de défoliation visant à tester la résilience des arbres, assuré l'entretien d'un dispositif de 1~000 arbres, installé des dendromètres et analysé les réponses de croissance aux traitements. 

% --- --- --- --- --- --- --- --- --- --- --- --- --- --- --- --- --- --- ---
% Item 5 %
% --- --- --- --- --- --- --- --- --- --- --- --- --- --- --- --- --- --- ---
\textbf{Bourse de recherche de premier cycle (BRPC/USRA), Université du Québec à Montréal} --- mai à août 2022. Direction~: P. Drapeau.

Travaux de terrain en Abitie sur le domaine vital du Grand Pic~: baguage, recherche de nids et classification des arbres. J'ai ensuite utilisé 20 ans de données de nidification pour tester les changements phénologiques de nidification de l'espèce en réponse aux changements climatiques~; il s'agissait de mon premier projet statistique mené de façon autonome. 

% --- --- --- --- --- --- --- --- --- --- --- --- --- --- --- --- --- --- ---
% Item 6 %
% --- --- --- --- --- --- --- --- --- --- --- --- --- --- --- --- --- --- ---
\textbf{Assistanat de recherche, Université du Québec à Montréal} --- 2021. Direction~: P. Drapeau.
Mon entrée en recherche en sciences naturelles~: étude des dommages causés par les pics aux infrastructures électriques à l'échelle du Québec. J'y ai appris la rigueur des protocoles de terrain, la gestion de bases de données et la programmation scientifique.

% <><><><><><><><><><><><><><><><><><><><><><><><><><><><><><><><><><><><><><><>
\textbf{b. Enseignement et encadrement} 
% <><><><><><><><><><><><><><><><><><><><><><><><><><><><><><><><><><><><><><><>
\textbf{Auxiliaire d'enseignement, FRST304 (physiologie végétale), UBC} --- 2025--2026. 
\textbf{Auxiliaire d'enseignement, FRST303 (écologie forestière), UBC} --- 2024--2025. 
Dans les deux cours pour deux sessions respectives, j'ai offert du soutien aux étudiants, préparé du matériel complémentaire pour aider les étudiantes et étudiants à s'approprier les concepts d'écologie et de physiologie végétale, participé au matériel d'enseignement et aidé à la préparation des examens.

\textbf{Supervision d'étudiantes et d'étudiants de premier cycle, UBC} --- 2024--2026. 
Au cours de mon parcours à UBC, j'ai formé et encadré 15 étudiant-e-s au premier cycle, dont deux qui poursuivent aujourd'hui des études graduées

\textbf{Bénévolat, tutorat et mentorat, UQÀM et UBC} --- depuis 2021. 
J'ai accompagné 12 étudiantes et étudiants internationaux durant leurs premiers mois au Canada, (Programme Allô, UQÀM) donné du tutorat en biochimie et en statistiques à 7 étudiants au premier cycle, et encadré 8 étudiants sur l'admission et le parcours pour se rendre aux études graduées dans le cadre de \textit{Biology Undergraduate Diversity in Research}.

% ==============================================================================
\subsection*{II --- Réalisations scientifiques}
% ==============================================================================

% <><><><><><><><><><><><><><><><><><><><><><><><><><><><><><><><><><><><><><><>
\textbf{a. Articles publiés ou acceptés dans des revues avec comité de lecture} 
% <><><><><><><><><><><><><><><><><><><><><><><><><><><><><><><><><><><><><><><>

% --- --- --- --- --- --- --- --- --- --- --- --- --- --- --- --- --- --- ---
% Item 1 %
% --- --- --- --- --- --- --- --- --- --- --- --- --- --- --- --- --- --- ---
Wolkovich, E. M., \textbf{Rouleau-Desrochers, C.}, García de Cortázar-Atauri, I., Walker, M. A., \& Lacombe, T. (2025). Building a more predictive model of terroir for the Anthropocene. \textit{Plants, People, Planet}, \textit{7}(6), 1611--1620. https://doi.org/10.1002/ppp3.70055 (travaux de baccalauréat). 

\textit{Rôle}~: j'ai développé les figures dirigé l'analyse cartographique globale ayant permis d'identifier les principales lacunes géographiques au niveau des données~; j'ai collaboré à la rédaction et révisé les sections méthodes et implications.

% <><><><><><><><><><><><><><><><><><><><><><><><><><><><><><><><><><><><><><><>
\textbf{b. Manuscrits soumis à des revues avec comité de lecture} 
% <><><><><><><><><><><><><><><><><><><><><><><><><><><><><><><><><><><><><><><>

% --- --- --- --- --- --- --- --- --- --- --- --- --- --- --- --- --- --- ---
% Item 1 %
% --- --- --- --- --- --- --- --- --- --- --- --- --- --- --- --- --- --- ---
\textbf{Rouleau-Desrochers, C.}, Van der Meersch, V., Buonaiuto, D. M., Chamberlain, C. J., Loughnan, D., Pederson, N., \& Wolkovich, E. M. (2026). \textit{Warmer and longer growing seasons increase growth for juvenile trees in a common garden} [manuscrit en révision, 28 p.]. \textit{Global Change Biology}. (travaux de maîtrise). 

\textit{Rôle}~: premier auteur. J'ai formulé les questions de recherche, récolté et analysé les données, dirigé la rédaction ainsi que l'interprétation des résultats et de leurs implications pour la régénération forestière et la réponse des arbres aux changements climatiques.

% <><><><><><><><><><><><><><><><><><><><><><><><><><><><><><><><><><><><><><><>
\textbf{c. Communications scientifiques avec comité de lecture} 
% <><><><><><><><><><><><><><><><><><><><><><><><><><><><><><><><><><><><><><><>

% --- --- --- --- --- --- --- --- --- --- --- --- --- --- --- --- --- --- ---
% Item 1 %
% --- --- --- --- --- --- --- --- --- --- --- --- --- --- --- --- --- --- ---
\textbf{Rouleau-Desrochers, C.}*, Van der Meersch, V., Pederson, N., \& Wolkovich, E. M. (Mai 2026). \textit{Evidence that growing season length and tree growth are decoupled in an urban arboretum} [communication orale]. Congrès de la Société canadienne d'écologie et d'évolution (SCEE/CSEE), Toronto Canada. Présentation orale complète, congrès d'envergure nationale (travaux de maîtrise). 

% --- --- --- --- --- --- --- --- --- --- --- --- --- --- --- --- --- --- ---
% Item 2 %
% --- --- --- --- --- --- --- --- --- --- --- --- --- --- --- --- --- --- ---
\textbf{Rouleau-Desrochers, C.}*, Pederson, N., \& Wolkovich, E. M. (Novembre 2025). \textit{How growing season length affects tree growth in a common garden study} [communication orale]. Colloque EcoEvo, Vancouver, C.-B., Canada. Congrès d'envergure régionale~; mention honorable (travaux de maîtrise).

\noindent {\footnotesize * Personne ayant présenté la communication.}

% <><><><><><><><><><><><><><><><><><><><><><><><><><><><><><><><><><><><><><><>
\textbf{d. Publications sans comité de lecture} 
% <><><><><><><><><><><><><><><><><><><><><><><><><><><><><><><><><><><><><><><>

% --- --- --- --- --- --- --- --- --- --- --- --- --- --- --- --- --- --- ---
% Item 1 %
% --- --- --- --- --- --- --- --- --- --- --- --- --- --- --- --- --- --- ---
\textbf{Rouleau-Desrochers, C.}, Patry, É., Gingras, G., Ucler, J., Vézina-Doré, M., Paillé, O., \& Moghaddam, Y. Y. (2024). Un été parti en fumée. \textit{Le Point Bio}, \textit{18}, 9--19 (travaux de baccalauréat). 

\textit{Rôle}~: premier auteur. J'ai mené les entrevues auprès des chercheuses et chercheurs et rédigé le manuscrit.

% ==============================================================================
\subsection*{III --- Autres réalisations d'ordre professionnel et social}
% ==============================================================================

% <><><><><><><><><><><><><><><><><><><><><><><><><><><><><><><><><><><><><><><>
\textbf{a. Bourses et distinctions} 
% <><><><><><><><><><><><><><><><><><><><><><><><><><><><><><><><><><><><><><><>
Maîtrise (2024--2026)~: 3 bourses obtenues par concours totalisant 28~750~\$, mention honorable pour ma communication à EcoEvo et désignation \textit{Honours} pour mon mémoire de maîtrise. 

Baccalauréat (2021--2024)~: Médaille académique d'argent du Gouverneur général, mention d'honneur, liste d'excellence du doyen, 13 bourses obtenues par concours totalisant 76~750~\$ et moyenne cumulative de 4,30/4,30. 


% <><><><><><><><><><><><><><><><><><><><><><><><><><><><><><><><><><><><><><><>
\textbf{b. Diffusion des sciences et dialogue science--société} 
% <><><><><><><><><><><><><><><><><><><><><><><><><><><><><><><><><><><><><><><>
Présentations de mes travaux à la communauté de science citoyenne Treespotters de l'Arboretum Arnold (2025 et 2026)~; animation d'une journée portes ouvertes au congrès de la SCEE (2024)~; reportage de BioChambers sur mes travaux en chambre de croissance (2024)~; entrevue lors de la cérémonie des bourses de l'UQAM (2023). 
Ma photographie animalière prolonge cet effort de diffusion~: lauréat du concours de photographie du Biodiversity Research Centre (UBC) 2026~;  mes images ont été publiées dans les \textit{Proceedings of the National Academy of Sciences} (2025), dans \textit{Le Rappel} (2025) et une photo figure sur la page couverture du \textit{Point Bio} (2024).

% <><><><><><><><><><><><><><><><><><><><><><><><><><><><><><><><><><><><><><><>
\textbf{c. Formation quantitative complémentaire} 
% <><><><><><><><><><><><><><><><><><><><><><><><><><><><><><><><><><><><><><><>
Durant ma première année de maîtrise, j'ai consolidé mes compétences quantitatives par un cours de biomathématiques, un séminaire de trois semaines sur les modèles bayésiens hiérarchiques et le cours \textit{Statistical Rethinking}, tout en raffinant mes modèles au sein d'un groupe de travail bayésien bimensuel.

% <><><><><><><><><><><><><><><><><><><><><><><><><><><><><><><><><><><><><><><>
\textbf{d. Parcours professionnel antérieur} 
% <><><><><><><><><><><><><><><><><><><><><><><><><><><><><><><><><><><><><><><>
Avant mon parcours universitaire actuel, j'ai obtenu des diplômes en restauration gastronomique et en sommellerie et travaillé près de dix ans dans cette industrie, où j'ai développé une solide rigueur et éthique de travail. La transition vers le milieu académique s'est amorcée pendant les six mois de travail dans un vignoble en Nouvelle-Zélande, où j'ai compris que mes liens avec la nature étaient plus profonds que je ne le croyais. Je me suis ensuite inscrit au certificat en écologie qui a précédé le baccalauréat en biologique à l'UQÀM que j'ai complété au printemps 2024.

\end{document}
